\chapter*{Resumen}
\addcontentsline{toc}{chapter}{Resumen}

Este Trabajo Fin de Grado estudia la detección de intrusiones en red como una decisión binaria, sensible al coste y realizada flujo a flujo: \texttt{PERMIT} o \texttt{BLOCK}. Se ha construido un pipeline reproducible que transforma CICIDS2017 en una observación canónica de 152 dimensiones ---76 características de flujo y 76 indicadores de presencia---, excluye campos susceptibles de fuga, ajusta el preprocesamiento únicamente con la partición de entrenamiento y entrena un agente \emph{Quantile Regression Deep Q-Network} (QRDQN). La recompensa penaliza con mayor intensidad permitir un ataque que bloquear tráfico benigno. Como el factor de descuento se fija en $\gamma=0$, el problema se interpreta expresamente como un bandido contextual y no como una defensa secuencial que modifica una red real.

La contribución metodológica incluye un perfil experimental congelado, semillas separadas para partición y modelo, una caché canónica sin escalar, contratos completos de artefactos y una campaña final diseñada para distinguir aprendizaje intradistribución de generalización. Esta campaña comprende una ejecución MAIN nueva, validación directa, intervalos \emph{bootstrap}, análisis de duplicados, control con etiquetas barajadas, partición completa por días, escalera de tamaños, sensibilidad a la semilla del modelo, cuatro escenarios retenidos dirigidos, comparaciones equivalentes con \emph{Random Forest} e inferencia offline sobre tráfico de laboratorio.

En el momento de cerrar este borrador de supervisión, la campaña final todavía no ha producido resultados científicos. Los artefactos históricos muestran que el pipeline aprende con una partición aleatoria y, a la vez, que los duplicados y el cambio de día obligan a limitar cualquier afirmación de generalización; se presentan únicamente como evidencia previa a la campaña. Por ello no se anticipa una conclusión numérica final. El resultado ya defendible es un método auditable para determinar, cuando existan los artefactos finales, qué comportamiento persiste al variar distribución, volumen de entrenamiento, semilla, escenario y familia de modelo. El sistema permanece restringido a experimentación e inferencia offline y no se considera preparado para despliegue operativo.

\vspace{1.2em}
\noindent\textbf{Palabras clave:} aprendizaje por refuerzo; detección de intrusiones en red; CICIDS2017; generalización; QRDQN; reproducibilidad.

\clearpage
